% TOP_PROBLEM: {"id": 20002519, "title": "An invariant certificate for caustic-forming Kerr rays", "category_id": 16, "category": {"id": 16, "name": "physics", "display_name": "Mathematical Physics", "description": "Problems at the intersection of mathematics and physics.", "slug": "physics", "order_index": 16, "created_at": "2026-07-31T15:26:25.671Z"}, "status": "partially_solved", "statusQualification": "The notion of stability and the admissible perturbation class require a definition in any claimed theorem. The original includes numerical plotting and physical interpretation alongside mathematical existence and classification questions. The catalogue title/status was paired with a machine-generated research attempt; that attempt is not peer-reviewed evidence of settlement. The statement covers the original source questions rather than only the title’s subcase."}
% FIRST_POSTED: 2026-10-11
% ENTRY_KIND: statement_only
% SUBMISSION: {"kind": "statement_only", "category": "top_problems", "problem_id": "20002519", "proposal_type": "established", "source": "American Institute of Mathematics, Gravitational lensing in the Kerr spacetime geometry, Questions 14–23. https://aimath.org/WWN/lensing/lensing.pdf", "status_evidence": "Inherited source-catalogue status partially_solved preserved. Background has machine-generated research claims explicitly marked not peer reviewed; these are not independent status evidence. Primary source accessed 11 October 2026; compared against the supplied catalogue excerpt. Historical wording and partial-result remarks were checked for scope; no comprehensive current-literature or proof-verification audit was performed.", "status_checked_on": "2026-10-11", "duplicate_check": "Full 15,773-record source registry checked for identical normalized statements; available site registry checked for identical normalized title/statement and canonical-ID collision. Canonical ID is absent from the site registry. No mathematical-equivalence literature audit performed.", "attribution": "Proposer not named in the AIM problem entry. Statement preparation: OpenAI GPT-6 Codex assistant, generation 11 October 2026; source-matching and notation editing by the assistant, no independent human mathematical verification."}
% SOURCE_STATEMENT: "(14) Plotting the caustics from conjugate points: Which photons (light rays) form the caustics? (15) What does the lightcone look like? (16) Plot higher order caustic sheets numerically. (17) Plot primary caustic sheet for observer close to Black hole. (18) How do caustics look for observers outside the equatorial plane? (19) How are the astroid caustics formed? (20) Is the birth process of the astroid caustic sheet for Kerr stable? (21) Clarify GR definitions of twist/shear/convergence in terms of lensing definition. (22) Stability of caustics with respect to metric perturbations? (23) Are all stable caustic sheets lightlike?"
% Preparation and final self-audit: CONTRIBUTING.md prompts C and D; source/audit artifacts retained in the local physics_statements_2026_10_11 batch.
% Imported from: physics_statements.tex; UnsolvedMath ID 20002519
% !TeX program = lualatex
\documentclass[11pt]{article}
\usepackage[margin=25mm]{geometry}
\usepackage{fontspec}
\setmainfont{Latin Modern Roman}
\usepackage{amsmath,amssymb,mathtools}
\usepackage{unicode-math}
\setmathfont{Latin Modern Math}
\usepackage{xurl,hyperref}
\hypersetup{colorlinks=true,urlcolor=blue}
\setcounter{secnumdepth}{0}
\setlength{\parindent}{0pt}
\setlength{\parskip}{5pt}
\setlength{\emergencystretch}{3em}
\newcommand{\sref}[2]{\hyperlink{source:#1:#2}{[S#2]}}
\newcommand{\cataloguescope}[1]{\paragraph{Recorded target.}#1}
\newcommand{\R}{\mathbb R}
\newcommand{\C}{\mathbb C}
\newcommand{\N}{\mathbb N}
\newcommand{\Z}{\mathbb Z}
\newcommand{\Q}{\mathbb Q}
\newcommand{\F}{\mathbb F}
\newcommand{\eps}{\varepsilon}
\begin{document}
% UnsolvedMath ID 20002519; source catalogue code AIM-PHYSICS-0027
\setcounter{section}{20002518}
\section{An invariant certificate for caustic-forming Kerr rays}\label{20002519}
{\small UnsolvedMath ID 20002519 · Mathematical Physics}
\subsection{Definitions and mathematical statement}

\paragraph{Shared notation.}$\mathbb N=\{1,2,\ldots\}$; the standard number fields have their usual meanings. All additional conventions are stated locally.


\paragraph{Definitions and notation.}A null geodesic models a light ray. A conjugate point occurs when a nontrivial Jacobi field vanishes at both endpoints, producing failure of local invertibility of the geodesic endpoint map. A light cone is the image of null directions from an event. Optical expansion, shear and twist describe the trace, symmetric traceless and antisymmetric parts of the transverse derivative of a null-ray congruence.

\paragraph{Catalogue wording retained for traceability.}
Retain source questions (14)--(23): identify which photons generate caustics via conjugate points; determine the light-cone geometry; plot higher-order caustic sheets; plot the primary sheet for an observer near the black hole; describe sheets for off-equatorial observers; explain formation of astroid caustics; determine stability of that birth process for Kerr; relate general-relativistic twist, shear and convergence to lensing definitions; determine stability under metric perturbations; and decide whether all stable caustic sheets are lightlike.
\paragraph{Scope and source qualification.}The notion of stability and the admissible perturbation class require a definition in any claimed theorem. The original includes numerical plotting and physical interpretation alongside mathematical existence and classification questions. The catalogue title/status was paired with a machine-generated research attempt; that attempt is not peer-reviewed evidence of settlement. The statement covers the original source questions rather than only the title’s subcase.
\subsection{Short English statement}
Identify the rays forming Kerr caustics, compute their sheets and study their geometry and stability.
\subsection{Sources}
\begin{itemize}
\item[\hypertarget{source:20002519:1}{S1}] ULAM.AI and the UnsolvedMath Contributors, UnsolvedMath, record 20002519 (AIM-PHYSICS-0027). Catalogue curation under CC BY 4.0; definitions and English exposition added locally.
\url{https://huggingface.co/datasets/ulamai/UnsolvedMath}.
\item[\hypertarget{source:20002519:2}{S2}] American Institute of Mathematics, Gravitational lensing in the Kerr spacetime geometry, Questions 14–23.
\url{https://aimath.org/WWN/lensing/lensing.pdf}.
\end{itemize}
\label{end:20002519}
\end{document}
